\GXNSFBodyText

\noindent\textbf{三年研究计划}\par

\noindent\textbf{第一年：数据规范与事件链构建}\par
完成文献综述、数据合规方案和事件编码手册；以佐佐木希公开资料建立演示事件链，完成多语言实体对齐与来源分级；在小样本上比较人工分段、隐马尔可夫模型和变点检测结果。年末将形成可复核的数据字典、标注指南和状态识别基线，并完成保留样本的一致性检验。

\noindent\textbf{第二年：模型构建与扰动分析}\par
完成文本、视觉元数据和平台行为特征的联合表示，建立平台校准的隐状态模型；构造候选扰动窗口，比较中断时间序列、匹配对照与合成控制的适用条件；开展指标权重、缺失机制和时间窗口敏感性分析。年末将完成核心模型、留出时间窗预测、消融实验和阶段性论文草稿。

\noindent\textbf{第三年：外部验证与组件适用性测试}\par
引入同期同构人物对照样本，检验完整模型的外部效度；以单一文化内容 IP 或单一区域品牌的固定观察窗构建广西—东盟小规模样本，评估多语言实体对齐、平台校准和扰动窗口组件的误差；汇总模型不支持迁移的案例、可发布的派生数据和复现实验脚本，完成结题报告与成果交流。

\begin{table}[htbp]
  \centering
  \fontsize{10}{14}\selectfont
  \caption{年度任务与考核节点（示例）}
  \begin{tabularx}{\textwidth}{p{0.14\textwidth}X X}
    \toprule
    年度 & 主要任务 & 阶段性考核节点 \\
    \midrule
    第一年 & 事件链、编码手册、基线模型 & 编码规则修订；保留样本一致性检验 \\
    第二年 & 多模态状态模型、扰动分析 & 留出时间窗预测；消融和窗口/权重敏感性分析 \\
    第三年 & 同构验证、组件测试、成果汇总 & 同构样本迁移结果；组件误差与适用边界 \\
    \bottomrule
  \end{tabularx}
  \label{tab:gx-annual-plan}
\end{table}

\noindent\textbf{学术交流与合作计划}\par

\noindent\textbf{交流安排}\par
每年围绕计算社会科学、复杂网络或数字文化传播参加一次国内学术会议，并在项目中期组织一次跨学科研讨。交流重点是编码规范、因果识别假设和模型适用边界。合作单位、访问人员和会议名称将在正式申请时依据实际安排填写。
